\section{Planner's Optimization Problem}
\label{sec:optimization}

\subsection{Problem Formulation}

The social planner solves \citep{kantorovich1965best,dantzig1963linear}:
%
\begin{equation}\label{eq:planning_problem}
\begin{aligned}
    \max_{\vec{C}_t,\,\vec{X}_t} \quad
        & \hat{W}_t(\vec{C}_t) = \sum_{i=1}^n \alpha^{agg}_{i,t}
          \,\frac{C_{i,t}^{1-\sigma_i}}{1-\sigma_i} \\
    \text{s.t.}\quad
        & (I - \bar{A})\vec{X}_t = \vec{F}_{t}
          \equiv \vec{C}_t + \vec{G}_{t}+ \Delta\vec{K}_{t}, \\
        & \vec{l}^{\top} \vec{X}_t \leq L,
     \quad B\vec{X}_t \leq \vec{K}_t, \\
\end{aligned}
\end{equation}

\subsection{Lagrangian Formulation}

The Lagrangian is:
%
\begin{equation}\label{eq:lagrangian}
\begin{aligned}
    \mathcal{L} = \hat{W}_t(\vec{C}_t)
        + \lambda_{L} (L - \vec{l} \cdot \vec{X}_t)
        \\+ \vec{\lambda}_{K} \cdot (\vec{K}_t - B\vec{X}_t)
     + \vec{\pi} \cdot
          \bigl[(I - \bar{A})\vec{X}_t - \vec{F}_{t}\bigr],
\end{aligned}
\end{equation}
%
where:
\begin{itemize}
    \item $\lambda_L \geq 0$: shadow price of labour;
    \item $\vec{\lambda}_K \in \mathbb{R}^n_+$: shadow prices of capital
          goods;
    \item $\vec{\pi} \in \mathbb{R}^n$: shadow prices of produced
          goods.
\end{itemize}

\subsection{First-Order Conditions}

\begin{theorem}[Optimality Conditions]
\label{thm:foc}
At an interior optimum, the following conditions hold:
\begin{align}
    \frac{\partial \hat{W}_t}{\partial C_{i,t}}
    = \alpha^{agg}_{i,t}\,C_{i,t}^{-\sigma_i}
 &= \pi_i,
    \label{eq:foc_consumption}\\[4pt]
    (I - \bar{A})^{\top} \vec{\pi}
        &= \lambda_L \vec{l} + B^{\top} \vec{\lambda}_{K}.
    \label{eq:foc_production}
\end{align}
\end{theorem}

\begin{proof}
Taking partial derivatives of $\mathcal{L}$ with respect to $C_{i,t}$
and setting equal to zero gives
\[
    \frac{\partial \mathcal{L}}{\partial C_{i,t}}
    = \alpha^{agg}_{i,t}\,C_{i,t}^{-\sigma_i} - \pi_i = 0,
\]
yielding~\eqref{eq:foc_consumption}. Differentiating with respect to
$\vec{X}_t$ yields~\eqref{eq:foc_production} by analogous reasoning.
\end{proof}

Since $\pi_i = \mu_t P_{i,t}$ (equation~\eqref{eq:price_proportional}),
the shadow price conditions reduce to:
%
\begin{equation}\label{eq:shadow_prices}
    \pi_i = \mu_t P_{i,t} = \alpha^{agg}_{i,t}\,C_{i,t}^{-\sigma_i},
\end{equation}
%
consistently with equation~\eqref{eq:crra_demand}. Solving for
$C_{i,t}$ recovers the aggregate Marshallian demand:
%
\begin{equation}
    C_{i,t} = \left(\frac{\alpha^{agg}_{i,t}}{\mu_t P_{i,t}}\right)^{1/\sigma_i}.
\end{equation}

We can find the aggregate first-order condition reducing this to an
aggregate constraint for the planner. Summing over households and using
the household-level first-order condition
$w_h\,\alpha_{i,t,h}\,C_{h,i,t}^{-\sigma_i} = \mu_t P_{i,t}$
(which gives $C_{h,i,t} = (w_h\,\alpha_{i,t,h}/(\mu_t P_{i,t}))^{1/\sigma_i}$),
aggregate demand satisfies:
%
\begin{equation}\label{aggregate consumer FOC}
    C_{i,t} = \sum_{h=1}^{n_h} C_{h,i,t}
    = 
      \sum_{h=1}^{n_h} \frac{(w_h\,\alpha_{i,t,h})^{1/\sigma_i}}{(\mu_t P_{i,t})^{1/\sigma_i}},
\end{equation}
%
which, upon defining the CES aggregate preference weight
$\alpha^{agg}_{i,t}$ as in equation~\eqref{eq:crra_demand}, is
equivalent to maximising the aggregate welfare function:
%
\begin{equation}
    \hat{W}_t(\vec{C}) = \sum_{i=1}^n \alpha^{agg}_{i,t}
      \,\frac{C_{i}^{1-\sigma_i}}{1-\sigma_i}.
\end{equation}

\begin{corollary}[Shadow Price Interpretation]
The shadow price $\pi_i = \mu_t P_{i,t}$ represents the marginal
social utility of good $i$, equating the marginal utility of
consumption to the value of the resource cost at shadow prices.
\end{corollary}

The welfare weights can be endogenously determined by the disposable income distribution between the households. Following the household-level optimum, the weight of household $h$ in the aggregate CES preference weight is proportional to its income share in supernumerary expenditure:
\begin{equation} \label{welfare weights}
    w_{h} = \frac{Y_{h}}{Y}
\end{equation}

This also shows the welfare weights are directly determined by the distribution matrix $\Phi$:
%
\begin{equation}
    \vec{w}_{h} = \frac{\Phi \vec{Y}_{f}}{Y}
\end{equation}
%
It must be noted this only holds true for the specific functional form
of the welfare function that has been assumed. In particular, the CRRA
structure with common $\sigma_i$ across households ensures that
household-level and aggregate demands are consistent: the aggregate
CES weight $\alpha^{agg}_{i,t}$ defined in
equation~\eqref{eq:crra_demand} correctly recovers aggregate demand
regardless of the income distribution, as long as welfare weights are
set by equation~\eqref{welfare weights}.
